% TOP_PROBLEM: {"id": 1412, "title": "Degree Diameter Problem", "category_id": 3, "category": {"id": 3, "name": "graph_theory", "display_name": "Graph Theory", "description": "Problems involving graphs, networks, and their properties.", "slug": "graph-theory", "order_index": 3, "created_at": "2026-07-31T15:26:25.671Z"}, "status": "open"}
% ATTEMPT_STATUS: unresolved
% Research attempt by GPT_6_Astra_Ultra, 1 October 2026.
% Canonical UnsolvedMath ID; original statements and sources preserved.
% FIRST_POSTED: 2026-10-01
% Imported from: unsolvedmath_additional_500.tex; UnsolvedMath ID 1412
% !TeX program = lualatex
% Compile twice: lualatex 1412.tex
% Self-contained: no external bibliography, images, or \input files.
% Numeric section labels and visible numbers are original UnsolvedMath IDs.
\documentclass[11pt,a4paper]{article}
\usepackage[margin=24mm,headheight=14pt]{geometry}
\usepackage{fontspec}
\setmainfont{Latin Modern Roman}
\setsansfont{Latin Modern Sans}
\setmonofont{Latin Modern Mono}
\usepackage{amsmath,amssymb,mathtools}
\usepackage{unicode-math}
\setmathfont{Latin Modern Math}
\usepackage{microtype}
\usepackage{enumitem,xurl,needspace}
\usepackage[dvipsnames]{xcolor}
\usepackage{hyperref}
\hypersetup{unicode=true,colorlinks=true,linkcolor=MidnightBlue,urlcolor=MidnightBlue,
 pdftitle={UnsolvedMath Problem 1412},pdfauthor={GPT 6 Astra Ultra},bookmarksnumbered=true}
\usepackage{bookmark,tocloft,titlesec,fancyhdr}
\setlength{\cftsecnumwidth}{4.2em}
\cftsetpnumwidth{2.5em}
\cftsetrmarg{4em}
\titleformat{\section}{\Large\bfseries\raggedright}{\thesection}{0.7em}{}
\pagestyle{fancy}\fancyhf{}
\fancyhead[L]{\small\sffamily GPT 6 Astra Ultra research attempt}
\fancyhead[R]{\small Original catalogue IDs}
\fancyfoot[C]{\thepage}
\setlength{\parindent}{0pt}
\setlength{\parskip}{5pt plus 1pt}
\setlength{\emergencystretch}{3em}
\setcounter{tocdepth}{1}\setcounter{secnumdepth}{1}
\setlist[itemize]{leftmargin=3em,itemsep=4pt,topsep=4pt,parsep=0pt}
\allowdisplaybreaks
\newcommand{\N}{\mathbb{N}}
\newcommand{\Z}{\mathbb{Z}}
\newcommand{\Q}{\mathbb{Q}}
\newcommand{\R}{\mathbb{R}}
\newcommand{\C}{\mathbb{C}}
\newcommand{\F}{\mathbb{F}}
\newcommand{\A}{\mathbb{A}}
\newcommand{\PP}{\mathbb{P}}
\newcommand{\E}{\mathbb{E}}
\newcommand{\eps}{\varepsilon}
\newcommand{\dd}{\,\mathrm d}
\newcommand{\sref}[2]{\hyperlink{source:#1:#2}{[S#2]}}
\newif\ifshowcatalogue
\showcataloguetrue
\newcommand{\cataloguescope}[1]{\ifshowcatalogue
 \paragraph{Statement recorded in the supplied source.}
 #1\fi}
\begin{document}
% UnsolvedMath ID 1412; source catalogue code GRAPH-025

\setcounter{section}{1411}

\section{The Degree--Diameter Extremal Problem}\label{1412}

{\small\sffamily UnsolvedMath ID 1412 · Graph Theory}

\subsection{Definitions and mathematical statement}

\paragraph{Shared notation.}Unless a section says otherwise, $\N=\{1,2,\ldots\}$,
$\Z$ is the ring of integers, $\Q,\R,\C$ are the rational, real, and complex
fields, $[N]=\{1,\ldots,\lfloor N\rfloor\}$, and $\log$ is the natural
logarithm. For real functions of a parameter tending to infinity,
$f\ll g$ and $f=O(g)$ mean $|f|\leq Cg$ eventually for some fixed $C$;
$f\gg g$ reverses this comparison; $f=o(g)$ means $f/g\to0$;
$f\sim g$ means $f/g\to1$; and $f\asymp g$ means both order bounds.
A subscript on $O$ or $\ll$ specifies allowed constant dependence.
The expression $n^{o(1)}$ in an upper bound means growth smaller than
$n^\epsilon$ for each fixed $\epsilon>0$. Local definitions govern any
exception, finite-field convention, or use of the same symbol for another
object. An asymptotic claim is not automatically an effective algorithm.



For integers $d,k\geq1$, let $N(d,k)$ be the maximum number of vertices of a finite connected simple graph of maximum degree at most $d$ and diameter at most $k$. Diameter is the maximum, over vertex pairs, of their shortest-path distance in edges. Determine $N(d,k)$ over the parameter range rather than restricting to regular or vertex-transitive graphs. \sref{1412}{1} \sref{1412}{2}

\paragraph{Statement recorded in the supplied source.}
For given maximum degree d and diameter k, what is the largest possible number of vertices in a graph? \sref{1412}{1}

\paragraph{Residual task reported by the archive.}

The embedded review (2026-08-17) records: Determine exact values for the many non-optimal table entries and improve general constructions or upper bounds away from the Moore cases. \sref{1412}{1}

\subsection{Short English statement}

How many vertices can a connected graph have when both the maximum degree and the greatest shortest-path distance are prescribed?

\subsection{Research attempt}

\paragraph{Target and interpretation of tabulated constructions.}
The quantity is maximized over all connected simple graphs, including
irregular and asymmetric graphs. The primary description of Comellas's
version-11 table [E1] distinguishes its largest known constructions
from entries proved optimal. This attempt derives universal bounds and
several exact cases; it does not treat every construction in a table as
an exact value of $N(d,k)$.

\paragraph{The breadth-first bound with all small cases separated.}
For $d\ge2$, fix a vertex $x$ and expose vertices by distance from it.
There is at most one vertex at distance zero, at most $d$ at distance
one, and at most $d(d-1)^{i-1}$ at distance $i\ge1$. Every discovered
vertex other than $x$ has at least one incident edge leading to an
earlier layer and hence at most $d-1$ edges available to discover new
vertices. Collisions only reduce the layer sizes. Since the diameter
is at most $k$,
\[
 N(d,k)\le M(d,k):=1+d\sum_{i=0}^{k-1}(d-1)^i.
\]
For $d>2$ this is $1+d((d-1)^k-1)/(d-2)$; retaining the sum
avoids a meaningless division at $d=2$. At $d=1$, a connected simple
graph has at most two vertices, and $K_2$ attains the bound for every
$k\ge1$, so $N(1,k)=2$.

At diameter one every distinct pair must be adjacent. Thus
$N(d,1)=d+1$, attained by $K_{d+1}$. At degree two, a connected
graph is a path or a cycle (or a single vertex). A path of diameter
at most $k$ has at most $k+1$ vertices, and a cycle of order $m$ has
diameter $\lfloor m/2\rfloor$. The cycle $C_{2k+1}$ therefore
attains the breadth-first bound:
\[
                        N(2,k)=2k+1.
\]
This also gives the correct $N(2,1)=3$.

\paragraph{An exact nontrivial case: $N(3,2)=10$.}
Use the ten two-element subsets of a five-element set as vertices,
joining two subsets when they are disjoint. A fixed subset has three
disjoint two-element subsets in its three-element complement, so the
graph is cubic. Two distinct nonadjacent vertices intersect in one
element; their union has three elements, and its two-element complement
is their unique common neighbour. Thus every pair is at distance at
most two. The graph has ten vertices, while $M(3,2)=1+3+6=10$.
Consequently it is optimal. This is the Petersen graph, derived here
from its subset model without assuming its extremality.

\paragraph{General constructions with controlled parameters.}
For an integer $q\ge2$, take words of length $k$ over an alphabet of
size $q$, joining words that differ in exactly one position. There
are $q^k$ vertices, the degree is $k(q-1)$, and the distance between
two words equals the number of positions in which they differ.
The upper bound on the distance follows by changing those positions
one at a time; the lower bound follows because one edge changes only
one position. Choosing $q=\lfloor d/k\rfloor+1$ gives
\[
                N(d,k)\ge (\lfloor d/k\rfloor+1)^k
                \quad\hbox{when }d\ge k.
\]
The graph has diameter exactly $k$. This is only a useful lower bound
in part of the parameter range, and the complete-graph lower bound
$N(d,k)\ge d+1$ remains available for every $k\ge1$.

A different word construction handles other ratios. For $q\ge2$,
join each word $(a_1,\ldots,a_k)$ to each word
$(a_2,\ldots,a_k,b)$ and make the edges undirected, deleting loops
and duplicate edges. There are at most $q$ successors and $q$
predecessors, so the maximum degree is at most $2q$. Starting with
any word, append the $k$ symbols of a target word in order; after
at most $k$ edge moves the target is reached. A repeated word in this
walk merely deletes a loop step. Thus for $d\ge4$,
\[
                       N(d,k)\ge\lfloor d/2\rfloor^k.
\]
Neither construction is asserted to give the largest currently known
graph for an individual pair $(d,k)$.

\paragraph{What equality would force.}
Attaining the breadth-first bound means every allowable child slot in
every layer is occupied by a distinct vertex. At $k=2$, this forces
every vertex to have degree $d$, no triangles, and no four-cycles.
For example two neighbours of the root cannot be joined, since that
wastes a child slot; nor can a second-layer vertex have two parents.
As every nonadjacent pair must then have exactly one common neighbour,
the adjacency matrix satisfies
\[
                         A^2=(d-1)I-A+J.
\]
This provides necessary algebraic restrictions on equality, but it
does not determine the best order when equality is impossible.
An upper bound excluding $M(d,k)$ leaves all smaller candidate orders
to be compared with actual constructions.

\paragraph{Executed checks and explicit residual task.}
The graph script constructs the subset model, checks every distance,
and obtains order ten, constant degree three, and diameter two. It also
checks its common-neighbour counts. The files are
\texttt{checks\_graphs.py} and \texttt{results\_graphs.json} in
\path{_Local/Erdos_problems/UnsolvedMath/attempt_audit_2026_10_01/number_theory/}.
No enumeration over all graph orders or all parameter pairs is claimed.
The exact answers above cover $d=1$, $d=2$, $k=1$, and $(d,k)=(3,2)$;
the general inequalities leave a substantial gap elsewhere. Determining
that gap requires either stronger universal exclusions or constructions
meeting them, not just a larger sample of symmetric graphs.
\textbf{Outcome: UNRESOLVED} for the catalogue-wide extremal problem.

\subsection{Sources}

{\small

\begin{itemize}

\item[\hypertarget{source:1412:1}{S1}] ULAM.AI, \emph{UnsolvedMath}, supplied \texttt{problems.json}, record 1412 (GRAPH-025), statement and background. The embedded literature-review date is 2026-08-17; that is the archive’s review date, not a fresh certification. Dataset landing page: \url{https://huggingface.co/datasets/ulamai/UnsolvedMath}.

\item[\hypertarget{source:1412:2}{S2}] F. Comellas, Table of Large Degree/Diameter Graphs, Mendeley Data, version 11 (January 2026). \url{https://doi.org/10.17632/d75dzbjd4k.11}. Bibliographic information imported from the supplied record.


\item[E1] Francesc Comellas, \emph{Table of Large
Degree/Diameter Graphs}, Mendeley Data, version 11, 28 January 2026.
\url{https://data.mendeley.com/datasets/d75dzbjd4k/11}.
Primary dataset description inspected 1 October 2026; it explicitly
distinguishes optimal entries from largest known constructions.

\end{itemize}

}

\label{end:1412}
\end{document}
